% GENERATED FILE. Edit canonical lessons, then run npm run generate:books.
% canonical-source-hash: fnv1a64:6e8b8a758bf35c62
% canonical-lessons: HI-C476-cehra, HI-C476-bhaumh, HI-C476-much, HI-C476-darhi, HI-C476-jabra

\chapter{The Face}
\label{ch:hi-the-face}
\hlchaptermodality{476}

\begin{chapteropening}
\textbf{By the end of this chapter:} \emph{I can describe someone's face.}

Complete the last lesson of chapter 476: I can describe someone's face.
\end{chapteropening}

\section[\texorpdfstring{\hi{चेहरा} (cehrā)}{चेहरा (cehrā)}]{\hi{चेहरा} (cehrā) — a face}

\label{lesson:HI-C476-cehra}



\begin{quote}
Before the new one: say the Hindi for a virtue, then the Hindi for a tendency.
\end{quote}

\subsection*{You'll want to know: \hi{चेहरा}}
\textbf{\hi{चेहरा}} — \emph{cehrā} — \textquotedblleft{}a face\textquotedblright{}.

\begin{quote}
\hi{धूप} \hi{में} \hi{मीरा} \hi{का} \hi{चेहरा} \hi{लाल} \hi{हो} \hi{गया।}
\end{quote}

\begin{quote}
\emph{dhūp meṁ Mīrā kā cehrā lāl ho gayā.}
\end{quote}

\begin{quote}
Mira's face went red in the sun.
\end{quote}

Masculine.

\begin{cousinweb}
\textbf{\hi{चेहरा}} (\emph{cehrā}) is usually traced to Persian \emph{chihra}, \textquotedblleft{}face, countenance\textquotedblright{}. \emph{mũh} is the face you wash; \emph{cehrā} is the face you see.
\end{cousinweb}

\begin{grammarlens}[title={What you've built}]
One more word for describing what you lived and felt.
\end{grammarlens}

\subsection*{Your turn}
\begin{itemize}
\raggedright
  \item \emph{Say it:} \emph{cehrā}
  \item \emph{Say it:} \emph{cehrā}, once more
  \item \emph{Say it:} \emph{dhūp meṁ Mīrā kā cehrā lāl ho gayā.} — Mira's face went red in the sun.
  \item \emph{Recall:} say the Hindi for a virtue, then the Hindi for a tendency, then say \emph{cehrā} again
\end{itemize}

\begin{tcolorbox}[breakable,colback=teal!4,colframe=teal!35!black,title={Before you move on}]
What does \hi{चेहरा} mean? (\textquotedblleft{}A face\textquotedblright{}.) Say it once more.
\end{tcolorbox}
\section[\texorpdfstring{\hi{भौंह} (bhauṁh)}{भौंह (bhauṁh)}]{\hi{भौंह} (bhauṁh) — an eyebrow}

\label{lesson:HI-C476-bhaumh}



\begin{quote}
Before the new one: say the Hindi for a tendency, then the Hindi for a face.
\end{quote}

\subsection*{You'll want to know: \hi{भौंह}}
\textbf{\hi{भौंह}} — \emph{bhauṁh} — \textquotedblleft{}an eyebrow\textquotedblright{}.

\begin{quote}
\hi{दादा} \hi{की} \hi{भौंहें} \hi{सफ़ेद} \hi{हैं।}
\end{quote}

\begin{quote}
\emph{dādā kī bhauṁheṁ safed haiṁ.}
\end{quote}

\begin{quote}
Grandfather's eyebrows are white.
\end{quote}

Feminine.

\begin{cousinweb}
\textbf{\hi{भौंह}} (\emph{bhauṁh}) is a tadbhava: CDIAL lists Hindi \emph{bhaũh}, \textquotedblleft{}eyebrow\textquotedblright{}, under Sanskrit \emph{bhrū}, \textquotedblleft{}eyebrow\textquotedblright{} (CDIAL 9688).
\end{cousinweb}

\begin{grammarlens}[title={What you've built}]
One more word for describing what you lived and felt.
\end{grammarlens}

\subsection*{Your turn}
\begin{itemize}
\raggedright
  \item \emph{Say it:} \emph{bhauṁh}
  \item \emph{Say it:} \emph{bhauṁh}, once more
  \item \emph{Say it:} \emph{dādā kī bhauṁheṁ safed haiṁ.} — Grandfather's eyebrows are white.
  \item \emph{Recall:} say the Hindi for a tendency, then the Hindi for a face, then say \emph{bhauṁh} again
\end{itemize}

\begin{tcolorbox}[breakable,colback=teal!4,colframe=teal!35!black,title={Before you move on}]
What does \hi{भौंह} mean? (\textquotedblleft{}An eyebrow\textquotedblright{}.) Say it once more.
\end{tcolorbox}
\section[\texorpdfstring{\hi{मूँछ} (mū̃ch)}{मूँछ (mū̃ch)}]{\hi{मूँछ} (mū̃ch) — a moustache}

\label{lesson:HI-C476-much}



\begin{quote}
Before the new one: say the Hindi for a face, then the Hindi for an eyebrow.
\end{quote}

\subsection*{You'll want to know: \hi{मूँछ}}
\textbf{\hi{मूँछ}} — \emph{mū̃ch} — \textquotedblleft{}a moustache\textquotedblright{}.

\begin{quote}
\hi{नाना} \hi{की} \hi{मूँछें} \hi{बहुत} \hi{लंबी} \hi{हैं।}
\end{quote}

\begin{quote}
\emph{nānā kī mū̃cheṁ bahut lambī haiṁ.}
\end{quote}

\begin{quote}
Grandfather has a very long moustache.
\end{quote}

Feminine.

\begin{cousinweb}
\textbf{\hi{मूँछ}} (\emph{mū̃ch}) is a tadbhava: CDIAL lists Hindi \emph{mūch}, \textquotedblleft{}moustache\textquotedblright{}, under \emph{*mucchā}, a form it reconstructs (CDIAL 10180). Hindi often uses the plural, \emph{mū̃cheṁ}, for one moustache.
\end{cousinweb}

\begin{grammarlens}[title={What you've built}]
One more word for describing what you lived and felt.
\end{grammarlens}

\subsection*{Your turn}
\begin{itemize}
\raggedright
  \item \emph{Say it:} \emph{mū̃ch}
  \item \emph{Say it:} \emph{mū̃ch}, once more
  \item \emph{Say it:} \emph{nānā kī mū̃cheṁ bahut lambī haiṁ.} — Grandfather has a very long moustache.
  \item \emph{Recall:} say the Hindi for a face, then the Hindi for an eyebrow, then say \emph{mū̃ch} again
\end{itemize}

\begin{tcolorbox}[breakable,colback=teal!4,colframe=teal!35!black,title={Before you move on}]
What does \hi{मूँछ} mean? (\textquotedblleft{}A moustache\textquotedblright{}.) Say it once more.
\end{tcolorbox}
\section[\texorpdfstring{\hi{दाढ़ी} (dāṛhī)}{दाढ़ी (dāṛhī)}]{\hi{दाढ़ी} (dāṛhī) — a beard}

\label{lesson:HI-C476-darhi}



\begin{quote}
Before the new one: say the Hindi for an eyebrow, then the Hindi for a moustache.
\end{quote}

\subsection*{You'll want to know: \hi{दाढ़ी}}
\textbf{\hi{दाढ़ी}} — \emph{dāṛhī} — \textquotedblleft{}a beard\textquotedblright{}.

\begin{quote}
\hi{चाचा} \hi{की} \hi{दाढ़ी} \hi{सफ़ेद} \hi{हो} \hi{गई} \hi{है।}
\end{quote}

\begin{quote}
\emph{cācā kī dāṛhī safed ho gaī hai.}
\end{quote}

\begin{quote}
My uncle's beard has gone white.
\end{quote}

Feminine.

\begin{cousinweb}
\textbf{\hi{दाढ़ी}} (\emph{dāṛhī}) is a tadbhava: CDIAL lists Hindi \emph{dāṛhī}, \textquotedblleft{}beard\textquotedblright{}, under \emph{*dāṁṣṭra}, a form it reconstructs beside Sanskrit \emph{daṁṣṭrā}, \textquotedblleft{}tusk, fang\textquotedblright{} (CDIAL 6250): the hair of the jaw.
\end{cousinweb}

\begin{grammarlens}[title={What you've built}]
One more word for describing what you lived and felt.
\end{grammarlens}

\subsection*{Your turn}
\begin{itemize}
\raggedright
  \item \emph{Say it:} \emph{dāṛhī}
  \item \emph{Say it:} \emph{dāṛhī}, once more
  \item \emph{Say it:} \emph{cācā kī dāṛhī safed ho gaī hai.} — My uncle's beard has gone white.
  \item \emph{Recall:} say the Hindi for an eyebrow, then the Hindi for a moustache, then say \emph{dāṛhī} again
\end{itemize}

\begin{tcolorbox}[breakable,colback=teal!4,colframe=teal!35!black,title={Before you move on}]
What does \hi{दाढ़ी} mean? (\textquotedblleft{}A beard\textquotedblright{}.) Say it once more.
\end{tcolorbox}
\section[\texorpdfstring{\hi{जबड़ा} (jabṛā)}{जबड़ा (jabṛā)}]{\hi{जबड़ा} (jabṛā) — a jaw}

\label{lesson:HI-C476-jabra}



\begin{quote}
Before the new one: say the Hindi for a moustache, then the Hindi for a beard.
\end{quote}

\subsection*{You'll want to know: \hi{जबड़ा}}
\textbf{\hi{जबड़ा}} — \emph{jabṛā} — \textquotedblleft{}a jaw\textquotedblright{}.

\begin{quote}
\hi{गिरने} \hi{से} \hi{अरुण} \hi{के} \hi{जबड़े} \hi{में} \hi{चोट} \hi{लगी।}
\end{quote}

\begin{quote}
\emph{girne se Aruṇ ke jabṛe meṁ coṭ lagī.}
\end{quote}

\begin{quote}
Arun hurt his jaw in a fall.
\end{quote}

Masculine.

\begin{cousinweb}
\textbf{\hi{जबड़ा}} (\emph{jabṛā}) is inherited: CDIAL lists Hindi \emph{jābṛā}, \textquotedblleft{}lower jaw-bone, jaw\textquotedblright{}, under Sanskrit \emph{jámbha}, \textquotedblleft{}tooth, jaw\textquotedblright{} (CDIAL 5137); \emph{jabṛā} is the common spoken form.
\end{cousinweb}

\begin{grammarlens}[title={What you've built}]
One more word for describing what you lived and felt.
\end{grammarlens}

\subsection*{Your turn}
\begin{itemize}
\raggedright
  \item \emph{Say it:} \emph{jabṛā}
  \item \emph{Say it:} \emph{jabṛā}, once more
  \item \emph{Say it:} \emph{girne se Aruṇ ke jabṛe meṁ coṭ lagī.} — Arun hurt his jaw in a fall.
  \item \emph{Recall:} say the Hindi for a moustache, then the Hindi for a beard, then say \emph{jabṛā} again
\end{itemize}

\begin{tcolorbox}[breakable,colback=teal!4,colframe=teal!35!black,title={Before you move on}]
What does \hi{जबड़ा} mean? (\textquotedblleft{}A jaw\textquotedblright{}.) Say it once more.
\end{tcolorbox}
